\section{Conclusion}

This paper compared two classical replanning baselines --- naive Dijkstra and heuristic-pruned A* --- against a Double DQN + HER policy trained with curriculum learning, for UAV routing in a shared dynamic grid environment. The evaluation used 50 matched scenarios with identical start-goal pairs and shared dynamic obstacles across all three methods. Dijkstra and A* both succeeded on all 50 scenarios and, as guaranteed by their shared optimality on this graph, produced identical path costs; RL succeeded on 49. On paired successful scenarios, both classical planners produced significantly shorter or equal paths than RL ($W=0$, $p=5.64\times10^{-7}$ vs.\ Dijkstra; $p=5.23\times10^{-7}$ vs.\ A*).

Adding A* changed the compute-time conclusion, not just extended it. Naive Dijkstra showed no statistically significant compute-time advantage over RL ($p=0.063$), which the original evaluation attributed to the grid being small enough that repeated full-graph replanning stays cheap. A* resolved the identical replanning problem roughly 16$\times$ faster than Dijkstra ($p=3.6\times10^{-15}$) by pruning its search with an admissible heuristic, while returning exactly the same path costs. RL was the slowest of the three methods in route-level compute time. The earlier explanation --- that RL's compute-time advantage might simply be latent at this scale --- was measuring RL against an unnecessarily naive classical opponent; a properly implemented classical planner is both more optimal and faster than RL here, not merely more optimal.

On the static 40-scenario evaluation, independently re-run for this revision, Dijkstra succeeded on every scenario while DQN succeeded on 97.5\%, with a mean path-cost penalty of 0.92 units above optimal and a re-measured $\sim$3.3$\times$ compute-time overhead for DQN on successful scenarios, consistent with the original study's static findings.

For the current controlled 15$\times$15 setup, classical replanning --- especially once heuristic pruning is used --- remains superior to the tested RL policy in reliability, path optimality, and compute cost. The RL policy demonstrates meaningful adaptability but does not outperform either classical baseline on any measured axis at this scale. This conclusion should be read alongside two explicit caveats this revision adds rather than resolves: all RL results come from a single trained policy (one random seed, no variance estimate), and the benchmark's small, clean scale was a deliberate methodological choice to maximize internal validity, not a claim about performance at operational UAV scale. Future work should test whether RL's balance of tradeoffs improves in larger, denser, or more frequently changing environments, against an incremental classical planner such as D* Lite, and across multiple RL training seeds.
